\documentclass[11pt]{article}
\usepackage[top=.50in, bottom=1.0in, left=1in, right=1in]{geometry}
\renewcommand{\baselinestretch}{1.1}
\usepackage{graphicx}
\usepackage{natbib}
\usepackage{amsmath}
\usepackage{amsfonts}
\bibliographystyle{besjournals}


\setlength\parindent{0pt}

\begin{document}

One tree, $\mathrm{N}$ ringwidth measurements $r_w$ for each year $x_n$.\\
For now, only a long-term Gaussian process:

\begin{equation}\label{eq:1}
\forall x_n  \hspace{0.5cm} \log{r_w(x_n)} \sim \text{normal}(\mu + f_{\text{long}}(x_n), \sigma^2)
\end{equation}

with:

\begin {equation}\label{eq:2}
\vec{f_{\text{long}}} \sim \text{multinormal}(\vec{0}, \mathrm{K_{\text{long}}})
\end{equation}

\vspace{0.5cm}
Each element $\mathrm{k}$ of the matrix $\underset{N\times N}{\mathrm{K}_{\text{long}}}$ is defined by the exponentiated quadratic function:
\begin {equation}
\forall i,j  \in [\![1,\mathrm{N}]\!] \hspace{1cm} \mathrm{k}_{i,j} = {\gamma_{\text{long}}^2} \exp\left(-\frac{1}{2}\left(\frac{\lvert x_j-x_i \rvert }{\rho_{\text{long}}}\right)^2\right)
\end{equation}
with $\gamma_{\text{long}}$ the marginal deviation and $\rho_{\text{long}}$ the length scale.

\vspace{0.5cm}
Note that equations (\ref{eq:1}) and (\ref{eq:2}) are equivalent to:

\begin{equation}
\log\vec{{r_w}} \sim \text{multinormal}(\vec{\mu}, \hspace{0.08cm} \Sigma + \mathrm{K}_{\text{long}})
\end{equation}

where:

\begin{equation}
    \vec{\mu} = \begin{bmatrix}
               \mu \\
               \mu \\
               \vdots \\
               \mu
    \end{bmatrix}
    \hspace{1cm}
	\Sigma = 
	\begin{bmatrix}
		\sigma^2  & 0 & \ldots & 0 \\
		0 & \ \sigma^2  & \ldots  & \vdots\\
		\vdots & \vdots & \ddots & \vdots \\
		0 & \ldots &  \ldots & \sigma^2 
	\end{bmatrix}
\end{equation}

\vspace{0.5cm}
For one tree, it's pretty straightforward to add a short-term Gaussian process:

\begin{equation}
\log\vec{{r_w}} \sim \text{multinormal}(\vec{\mu}, \hspace{0.08cm} \Sigma + \mathrm{K}_{\text{long}} + \mathrm{K}_{\text{short}})
\end{equation}

This short-term Gaussian process is also defined by an exponentiated quadratic function with parameters $\gamma_{\text{short}}$ and $\rho_{\text{short}}$.
\vspace{0.5cm}

In this case, you can write only one Gaussian process, with one matrix $\mathrm{K}_{\text{long+short}}$ with two exponentiated quadratic functions:
\begin {equation}
\forall i,j  \in [\![1,\mathrm{N}]\!] \hspace{1cm} \mathrm{k}_{i,j} = {\gamma_{\text{long}}^2} \exp\left(-\frac{1}{2}\left(\frac{\lvert x_j-x_i \rvert }{\rho_{\text{long}}}\right)^2\right) + {\gamma_{\text{short}}^2} \exp\left(-\frac{1}{2}\left(\frac{\lvert x_j-x_i \rvert }{\rho_{\text{short}}}\right)^2\right)
\end{equation}

\vspace{0.8cm}
Several trees $t \in [\![1,\mathrm{T}]\!]$ on the same stand, with the same number of ringwidth measurements:

\begin{equation}
\forall t  \in [\![1,\mathrm{T}]\!] \hspace{1cm} \log\vec{{r_w({t})}} \sim \text{multinormal}(\vec{\mu}, \hspace{0.08cm} \Sigma + \mathrm{K}_{\text{long+short}})
\end{equation}



\end{document}